\documentclass[11pt,a4paper]{article}
\usepackage[margin=25mm]{geometry}
\usepackage{amsmath,amssymb,amsthm,hyperref}
\newtheorem{lemma}{Lemma}
\title{An obstruction to an intermediate Suzuki classification}
\author{Human Compression project: F003 proof audit}
\date{}
\begin{document}\maketitle
This note concerns the intermediate classification in Section~3 of
P. Spiga, \emph{On the rank of Suzuki polytopes: an answer to Hubard
and Leemans}, arXiv:2010.04080v1, 8 October 2020. It does not prove
or disprove that paper's theorem about chiral polytopes. The inspected
four-page preprint is the fixed source; no journal body is substituted.

\begin{lemma}
Let $G$ be any group and let $\iota\in G$ satisfy $\iota^2=1$.
Put $X_\iota=\{x\in G:x\iota x=\iota\}$. Then
\[
 X_\iota\longrightarrow\{y\in G:y^2=1\},\qquad x\longmapsto\iota x
\]
is a bijection, with inverse $y\mapsto\iota y$.
\end{lemma}
\begin{proof}
Multiplication on the left by $\iota$ turns $x\iota x=\iota$
into $(\iota x)^2=1$. Both implications hold, and left multiplication
by $\iota$ is its own inverse.
\end{proof}

The standard Suzuki structure already retained in the project's
accepted Suzuki application gives $q^2+1$ ovoid points, $q-1$
nonidentity central root involutions fixing each point, and a unique
fixed point for every involution. Consequently
\[
 \#\{g\in\operatorname{Sz}(q):|g|=2\}=(q^2+1)(q-1),\qquad
 |X_\iota|=1+(q^2+1)(q-1)=q^3-q^2+q.
\]
These are retained classical group-structure inputs, not a new
classification of Suzuki groups. The bijection itself is elementary
and needs none of them.

The source defines exactly this set $X$ in equation~(1), with its
displayed involution $\iota$. At the end of Section~3.6 it identifies
$X$ with $\{\iota\}$ together with a diagonal family indexed by
$\mathbb F_q^*$. That union has $q$ elements: the diagonal family
has $q-1$ distinct elements and the displayed antidiagonal involution
is not diagonal. For every allowed $q\ge8$ the two counts disagree.
The claimed intermediate classification is therefore false.

There is also a shorter involution obstruction. The central root
involutions commuting with $\iota$ all belong to $X_\iota$; there
are $q-1$ of them. The displayed diagonal family has odd order and
contains no nonidentity involution, so its union with $\{\iota\}$
has only one. This argument is independent of any convention for
the diagonal exponents.

The exact bounded diagnostic in \texttt{symmetric\_suzuki\_probe.g}
uses GAP's standard matrix model of $\operatorname{Sz}(8)$ and its
antidiagonal involution. It checks the bijection on the actual sets.
It obtains group order $29120$, $455$ nonidentity involutions,
$|X_\iota|=456$, involution centralizer order $64$, and seven
nonidentity commuting involutions. It also supplies a symmetric
commuting involution different from $\iota$. The literal matrix
model need not have the source's diagonal convention: the
cardinality contradiction and the commuting-involution obstruction
are invariant under changing the Suzuki representation or conjugating
the chosen involution. No comparison of matrix entries across two
unproved model identifications is used.

\section*{What this leaves to prove}
The source states Wilson's product characterization only for
orthogonal vectors $u,v$. Its subsequent calculation also imposes
the displayed product identities on the basis pairs $(e_1,e_4)$
and $(e_2,e_3)$, which have symplectic pairing one. Their use is
not justified by that stated characterization. The false conclusion
above shows that the resulting system cannot classify all of $X$.
Fixing a notation or exponent alone does not repair the cardinality
contradiction.

The rank-three theorem could still be true. Its original chiral
polytope hypotheses include additional geometric and intersection
conditions that the asserted classification of all of $X$ does
not use. A replacement argument must exploit valid conditions of
that kind, or supply another proof. No rank-four polytope is
constructed here, and no named open question is claimed solved.
F003 remains pending complete recovery and independent review.
\end{document}
